\section{Related Work}
\label{sec:related_work}

We complement the cost-model background of \Cref{subsec:cost-models} with a
discussion of five threads most directly relevant to this paper:
cost-aware sampling for configurable systems, learning configuration spaces
in general, cost-aware sampling and active learning in machine learning,
input/workload sensitivity, and multi-objective evaluation.

\paragraph*{Cost-aware sampling for configurable systems.}
The closest prior work is Sarkar \etal~\cite{sarkar2015}, which introduces
progressive and projective sampling under a composite cost combining training
acquisition, testing acquisition, model building, and prediction error. As
discussed in \Cref{subsec:cost-models}, that paper collapses acquisition cost
to $2n$, baking in equal cost per configuration. Our work tests this
simplification empirically rather than adopting or rejecting it
\emph{a priori}. Westermann \etal~\cite{westermann2012} formulate efficiency
as ``required number of measurement points'' for a target error---another
count-based abstraction. The data-mining roots of these formulations are
Weiss and Tian's utility-based sampling and Last's projective sampling
\red{[cite Weiss \& Tian, ICDM/PKDD 2008; Last, KDD 2009 -- bib entries
needed]}, both of which fit partial learning curves and runtime curves
to estimate a utility-optimal training-set size. They reason in terms of
training-set size and aggregate runtime curves rather than per-configuration
heterogeneous measurement time.

\paragraph*{Learning configuration spaces.}
Beyond cost-aware sampling, a large body of work studies how to predict
performance of configurable systems with smaller and better-chosen
samples. Guo \etal~\cite{guo2013ASE,Guo2017} train predictors from samples of
size $N,2N,3N,M$. Siegmund
\etal~\cite{siegmund2012a,siegmund2013,siegmund2015FSE} build
performance-influence models with $t$-wise and option-frequency designs.
Nair \etal~\cite{nair2017,nair2018c,nair2018a} explore spectral, rank-based,
and progressive sampling. Jamshidi \etal~\cite{jamshidi2017b,jamshidi2018}
study transfer learning across environments. Temple
\etal~\cite{temple2017a} approach the problem from a learning-theoretic
angle. Kaltenecker \etal~\cite{kaltenecker2019} systematically compare
distance-based sampling to $t$-wise, solver-based, and randomized
solver-based. The systematic review of Pereira \etal~\cite{julianasurvey}
catalogues this lineage. Our work does not propose a new prediction approach
or a new sampler; it asks whether the cost axis under which these approaches
are compared is calibrated.

\paragraph*{Cost-aware sampling and active learning in ML.}
Outside configurable systems, the machine-learning community has
investigated heterogeneous per-instance acquisition cost for nearly two
decades. Settles, Craven, and Friedland's
\red{[cite Settles--Craven--Friedland, NIPS workshop 2008]} study of
\emph{active learning with real annotation costs} measures human labeling
times in four real-world domains and shows directly that
\emph{reducing the number of labeled instances does not necessarily reduce
total acquisition cost}---the same observation we make for configurable
systems, in a different domain. Beygelzimer \etal{}'s COAL
\red{[cite Beygelzimer et al., COAL]}, Margineantu's active cost-sensitive
learning \red{[cite Margineantu]}, and Arora \etal{}'s annotation-cost
prediction \red{[cite Arora--Nyberg--Rose]} all treat acquisition cost as a
heterogeneous, possibly predictable, per-instance quantity. The recent
survey of Herde \etal{} \red{[cite Herde et al.]} catalogues uniform and
heterogeneous cost types in active learning. Contardo
\etal{} \red{[cite Contardo et al.]} extend the picture to feature-acquisition
cost, where each feature has its own price.

The configurable-systems performance-prediction literature has not absorbed
this lesson: it has continued to use count-based proxies even though ML
active learning has known for years that count is not cost when instances
are heterogeneous. Our paper is the configurable-systems analogue of
Settles \etal: we measure heterogeneous per-configuration cost, test
whether sample count is a calibrated proxy, and ask whether replacing
counts with measured times changes scientific conclusions.

\paragraph*{Recent configurable-systems cost conventions.}
Recent papers by Tao Chen and collaborators provide an important nearby
comparison point because they also reason about expensive measurements,
tuning budgets, and resource efficiency for configurable systems. Tuning,
transfer, and prediction papers frequently express effort as numbers of
measurements, evaluations, or target-environment samples, because counts are
simple to compare, reproduce, and re-use across studies
\cite{chen2021multiobj,chen2023aspirations,chen2024admmo,gong2024sempl}. Some
studies already strengthen this practice by calibrating counts from measured
wall-clock budgets \cite{chen2021multiobj,chen2023aspirations}. Our
contribution is not to revisit those algorithms, but to isolate the common
untested assumption beneath them: whether count is a calibrated proxy for
measured acquisition cost, and whether conclusions survive when that assumption
is checked.

\paragraph*{Cost-aware regression vs. cost-aware acquisition.}
A reviewer might object that ``cost-aware regression'' is already a
literature, citing utility-based regression by Torgo and Ribeiro
\red{[cite Torgo--Ribeiro]}. That literature is conceptually distinct: it
weights \emph{prediction errors} by a domain-defined relevance/utility
function over the target value, so high-stakes errors count more than
low-stakes ones. It is about the asymmetric value of being wrong, not
about the cost of acquiring a label. Our work is about the latter. Both
notions of ``cost'' are legitimate, and neither subsumes the other.

\paragraph*{Model-building cost vs.\ acquisition cost.}
There is also recent Chen-line work on the cost of building configurable-system
performance models. Gong and Chen compare encoding schemes using both
prediction accuracy and model training time, including hyperparameter search
\cite{gong2022encoding}; Xiang \etal{} study measured online adaptation
overhead \cite{xiang2026dhda}; the recent DCPL survey maps the broader
deep-learning literature \cite{gong2024dcpl}; and Domland analyzes
context-sensitive landscape structure for tuning \cite{ye2026domland}. These
are real and relevant cost or context questions, but they are not the
acquisition-cost question we study here. We focus on the resource spent to
acquire labels for sampled configurations, not the compute spent fitting or
updating the predictor once those labels are available.

\paragraph*{Input and workload sensitivity.}
Pereira \etal~\cite{pereira2020} show that on x264, the ranking of sampling
strategies varies markedly across $17$ input videos and two performance
properties. Recent tuning-analysis work likewise argues that workload effects
are system-dependent rather than uniform \cite{ye2026domland}.
\red{[cite Lesoil et al., JSS 2023; Lesoil et al., JSS 2024;
Mühlbauer et al., ICSE 2023]} extend this point to broader collections of
configurable systems and to input-aware models, and the \dsWLI
study explicitly catalogues per-configuration measurement times across
$6$--$13$ workloads on each of nine systems---data we reuse in
\Cref{subsec:datasets}. Our use of multiple inputs and properties is not an
attempt to (re)establish that context matters. It is established. What is
new is the question whether the validity of the size-as-cost abstraction is
itself context-dependent: in some inputs $n$ may be a fine proxy, and in
others not.

\paragraph*{Multi-objective evaluation.}
Several works compare strategies on combined accuracy/effort
objectives. Sarkar \etal~\cite{sarkar2015} use a scalar composite. Siegmund
\etal~\cite{siegmund2015FSE} report Pareto trade-offs between accuracy and
number of measurements. Our composite objectives in
\Cref{subsec:composite} use the same shape but are evaluated under both
cost models, and are paired with a Pareto comparison
(\Cref{subsec:pareto},\Cref{subsec:rq4}) so the conclusions are not hostage
to a single $\lambda$. Li \etal~\cite{li2020pareto} further argue that
Pareto-based SBSE evaluation should select indicators that reflect the problem
nature and decision-maker preferences. Our use of frontier membership change
and overlap follows that guidance: we compare the same finite set of
strategy--budget candidates under two cost models, so the key question is
choice-set stability rather than generic front-approximation quality.

\paragraph*{What we are not.}
We are not proposing a new sampling strategy and we make no claim of
methodological novelty in cost-aware sampling \emph{in general}; the ML
active-learning lineage above already developed the heterogeneous-cost
framing. We are revisiting the cost term in the configurable-systems
performance-prediction literature, where the count-based abstraction has
remained the operational default. The closest empirical precedent in this
literature is Siegmund \etal~\cite{siegmund2012a}, which reports both the
number of measurements and the time in hours of each heuristic but treats
time as descriptive effort rather than as a proxy-validity question. We
frame the proxy-validity question explicitly, instrument it across systems,
inputs, properties and strategies, and compare conclusions under
count-based and empirical-cost objectives.
